\section{El assembler}
\label{sec:assembler}

El assembler es el programa \code{asm/pochoasm.py}, escrito en Python sin bibliotecas
externas (unas 400 líneas, incluidos los comentarios). Recibe un archivo de assembly y escribe un archivo
\code{.hex} con una palabra de 32 bits en hexadecimal por línea. La palabra de la línea
$k$ queda en la dirección $4k$ de la RAM, que es lo que espera \code{\$readmemh} en
\code{pochoco\_ram.v}.

\begin{lstlisting}[language=bash, caption={Uso del assembler.}, label={lst:uso}]
python3 asm/pochoasm.py sw/game.s -o sw/game.hex [-l sw/game.lst] [-D NOMBRE=VALOR]
\end{lstlisting}

La opción \code{-l} escribe además un listado con la dirección, la palabra y la línea de
origen de cada instrucción (anexo~\ref{anx:listado}). La opción \code{-D} reemplaza el
valor de una constante \code{.equ} sin editar el archivo; se usa para simular el juego
con tiempos cortos (sección~\ref{sec:verificacion}).

\subsection{Sintaxis aceptada}

Cada línea tiene la forma
\[
  [\,\textit{etiqueta}\texttt{:}\,]~~[\,\textit{instrucción}~~\textit{operando}\texttt{,}~\textit{operando}\ldots\,]~~[\,\texttt{\#}~\textit{comentario}\,]
\]
y cualquiera de las tres partes puede faltar. Una línea puede tener varias etiquetas.

\begin{itemize}
  \item \textbf{Registros}: \code{x0} a \code{x15} o sus nombres ABI (\code{zero},
        \code{ra}, \code{sp}, \code{gp}, \code{tp}, \code{t0}--\code{t2},
        \code{s0}/\code{fp}, \code{s1}, \code{a0}--\code{a5}). Los registros que
        existen en RV32I pero no en RV32E (\code{x16}--\code{x31}, \code{a6},
        \code{s2}, \code{t3}, \ldots) producen un error explicativo.
  \item \textbf{Números}: decimales, hexadecimales (\code{0x}), binarios (\code{0b}),
        con signo opcional. También se acepta el nombre de una constante \code{.equ},
        con signo menos opcional.
  \item \textbf{Acceso a memoria}: \code{desplazamiento(registro)}, por ejemplo
        \code{lw t0, CYCLES(s0)}. Si el desplazamiento se omite vale 0.
  \item \textbf{Destinos de salto}: siempre una etiqueta.
  \item \textbf{Instrucciones}: las 40 de RV32I: las diez de registro-registro
        (\code{add}, \code{sub}, \code{sll}, \code{slt}, \code{sltu}, \code{xor},
        \code{srl}, \code{sra}, \code{or}, \code{and}), las nueve con inmediato
        (\code{addi}, \code{slti}, \code{sltiu}, \code{xori}, \code{ori}, \code{andi},
        \code{slli}, \code{srli}, \code{srai}), cinco cargas, tres almacenamientos,
        seis branches, \code{lui}, \code{auipc}, \code{jal}, \code{jalr}, \code{ecall},
        \code{ebreak} y \code{fence}.
  \item \textbf{Directivas}: \code{.equ NOMBRE, valor} define una constante. Las
        directivas \code{.text}, \code{.section}, \code{.global} y \code{.globl} se
        aceptan y se ignoran: sirven al assembler de GNU para enlazar varios archivos,
        y aquí hay un solo programa en una sola memoria desde la dirección 0.
\end{itemize}

Las pseudo-instrucciones son abreviaturas estándar de RISC-V \cite{riscv-unpriv} que
no existen en el hardware. El assembler las reescribe como instrucciones reales
(tabla~\ref{tab:pseudo}).

\begin{table}[H]
\centering
\caption{Pseudo-instrucciones aceptadas y su traducción.}
\label{tab:pseudo}
\begin{tabular}{ll}
\toprule
Pseudo-instrucción & Instrucción real \\
\midrule
\code{nop} & \code{addi zero, zero, 0} \\
\code{mv rd, rs} & \code{addi rd, rs, 0} \\
\code{not rd, rs} & \code{xori rd, rs, -1} \\
\code{neg rd, rs} & \code{sub rd, zero, rs} \\
\code{li rd, valor} & \code{addi}, \code{lui} o \code{lui} + \code{addi} (sección~\ref{sec:li}) \\
\code{j etiq} & \code{jal zero, etiq} \\
\code{jal etiq}, \code{call etiq} & \code{jal ra, etiq} \\
\code{jr rs} & \code{jalr zero, 0(rs)} \\
\code{ret} & \code{jalr zero, 0(ra)} \\
\code{beqz rs, etiq} / \code{bnez rs, etiq} & \code{beq} / \code{bne rs, zero, etiq} \\
\code{bgt a, b, etiq} / \code{ble a, b, etiq} & \code{blt} / \code{bge b, a, etiq} \\
\code{bgtu a, b, etiq} / \code{bleu a, b, etiq} & \code{bltu} / \code{bgeu b, a, etiq} \\
\bottomrule
\end{tabular}
\end{table}

En el assembler de GNU, \code{call} se traduce como \code{auipc} + \code{jalr} para
alcanzar cualquier dirección de 32 bits. Aquí se traduce como un solo \code{jal}, cuyo
alcance de $\pm 1$~MB sobra para una RAM de 2~KB.

\subsection{Funcionamiento: tres pasadas}

El problema central de un assembler son las etiquetas usadas antes de definirse: en
\code{call show\_number}, la subrutina está al final del archivo y todavía no se sabe
en qué dirección quedará. Por eso el texto se recorre tres veces:

\begin{enumerate}
  \setcounter{enumi}{-1}
  \item \textbf{Constantes.} Se separa cada línea en etiquetas, mnemónico y operandos, y
        se evalúan las \code{.equ} en orden.
  \item \textbf{Direcciones.} Se recorre el programa con un contador de posición que
        parte en 0. Cada etiqueta recibe el valor actual del contador, y cada
        instrucción lo avanza en 4 bytes. Las pseudo-instrucciones se expanden aquí
        porque determinan cuántas palabras ocupa cada línea: \code{li} puede ocupar una
        o dos.
  \item \textbf{Codificación.} Con todas las direcciones conocidas, cada instrucción
        real se traduce a sus 32 bits.
\end{enumerate}

El tamaño de \code{li} se decide en la primera pasada y depende del valor, así que
\code{li} solo acepta números y constantes \code{.equ}, cuyo valor ya se conoce en ese
momento; no acepta etiquetas.

\subsection{Codificación}

RISC-V tiene seis formatos de instrucción \cite{riscv-unpriv, patterson-riscv}. Todos
dejan el código de operación en los bits 6:0 y los registros siempre en las mismas
posiciones: \code{rd} en 11:7, \code{rs1} en 19:15 y \code{rs2} en 24:20. Así el
hardware puede leer los registros antes de saber qué instrucción es. Lo que cambia es
dónde se reparte el inmediato (tabla~\ref{tab:formatos}).

\begin{table}[H]
\centering
\caption{Formatos de instrucción de RV32I, del bit 31 al 0.}
\label{tab:formatos}
\small
\begin{tabular}{l|c|c|c|c|c|c|}
\cline{2-7}
 & 31 \ldots 25 & 24 \ldots 20 & 19 \ldots 15 & 14 \ldots 12 & 11 \ldots 7 & 6 \ldots 0 \\
\cline{2-7}
R & funct7 & rs2 & rs1 & funct3 & rd & opcode \\
\cline{2-7}
I & \multicolumn{2}{c|}{imm[11:0]} & rs1 & funct3 & rd & opcode \\
\cline{2-7}
S & imm[11:5] & rs2 & rs1 & funct3 & imm[4:0] & opcode \\
\cline{2-7}
B & imm[12\textbar10:5] & rs2 & rs1 & funct3 & imm[4:1\textbar11] & opcode \\
\cline{2-7}
U & \multicolumn{4}{c|}{imm[31:12]} & rd & opcode \\
\cline{2-7}
J & \multicolumn{4}{c|}{imm[20\textbar10:1\textbar11\textbar19:12]} & rd & opcode \\
\cline{2-7}
\end{tabular}
\end{table}

En el assembler cada formato es una función de una línea que coloca cada campo con
desplazamientos y máscaras (código~\ref{lst:enc}). Los formatos B y J no guardan el
bit 0 del desplazamiento, porque las instrucciones están alineadas y siempre es 0. Por
eso duplican el alcance.

\begin{lstlisting}[language=Python, caption={Codificación de los formatos I y B en \code{pochoasm.py}.}, label={lst:enc}]
def enc_i(imm, rs1, f3, rd, op):
    return (imm & 0xFFF) << 20 | rs1 << 15 | f3 << 12 | rd << 7 | op

def enc_b(off, rs2, rs1, f3, op):
    # imm[12|10:5] en 31:25, imm[4:1|11] en 11:7 (el bit 0 no se guarda: siempre es 0)
    return (((off >> 12) & 1) << 31 | ((off >> 5) & 0x3F) << 25 | rs2 << 20 | rs1 << 15
            | f3 << 12 | ((off >> 1) & 0xF) << 8 | ((off >> 11) & 1) << 7 | op)
\end{lstlisting}

La información de cada instrucción está en una tabla de datos, no en código: el
mnemónico apunta a su formato, \code{opcode}, \code{funct3} y \code{funct7}. Por ejemplo,
\code{'bne': ('B', 0x63, 1, 0)}. Agregar una instrucción es agregar una fila.

\subsection{Etiquetas e inmediatos}

\paragraph{Etiquetas.} Los branches y \code{jal} no guardan la dirección de destino
sino la distancia desde la instrucción que salta: $\text{desplazamiento} =
\text{dirección de la etiqueta} - pc$. El procesador suma ese valor a su \code{pc}. Un
desplazamiento negativo es un salto hacia atrás, como el de un bucle.

\paragraph{Rangos.} Cada inmediato se verifica contra lo que cabe en su campo, con
signo: $[-2048, 2047]$ para los formatos I y S, $[0, 31]$ para la cantidad de
desplazamiento, $[0, 2^{20}-1]$ para \code{lui} y \code{auipc}, $[-4096, 4094]$ para
los branches y $[-2^{20}, 2^{20}-2]$ para \code{jal}.

\paragraph{La pseudo-instrucción \code{li}.}
\label{sec:li}
Carga un valor de 32 bits en un registro. Si cabe en 12 bits con signo, basta un
\code{addi} desde \code{zero}. Si no, se parte en una parte alta de 20 bits, que carga
\code{lui}, y una baja de 12 bits, que suma \code{addi}. La dificultad es que
\code{addi} interpreta sus 12 bits \emph{con signo}. Si el bit 11 de la parte baja
es 1, \code{addi} resta en vez de sumar, y la parte alta tiene que compensarlo sumando 1.
El assembler calcula:

\[
  \text{lo} = \text{extensión de signo de } (v \mathbin{\&} \texttt{0xFFF}), \qquad
  \text{hi} = \frac{v - \text{lo}}{2^{12}} \bmod 2^{20}.
\]

Si $\text{lo} = 0$ basta el \code{lui}. Este caso se verificó ejecutando en Python el
par \code{lui}/\code{addi} generado para 2.017 valores, incluidos los bordes
(\code{0x7FF}, \code{0x800}, \code{0xFFF}, \code{0x80000000}, \code{0xFFFFFFFF}) y
2.000 al azar.

\subsection{Validación}

Cada error se reporta con el archivo, la línea y el texto de la línea, y el assembler
sigue revisando el resto del archivo para mostrar todos los errores de una vez. Si hay
algún error no escribe el \code{.hex}, para que nunca quede en disco un archivo a medio
generar. El código~\ref{lst:errores} muestra un ejemplo.

\begin{lstlisting}[caption={Salida del assembler ante un programa con errores.}, label={lst:errores}]
malo.s:3: error: 'x16' no existe en RV32E: solo hay 16 registros (x0-x15)
    add  x16, a0, a1
malo.s:7: error: inmediato de 12 bits fuera de rango: 3000 (permitido -2048 a 2047)
    addi t0, t0, 3000
malo.s:9: error: etiqueta no definida: 'fin'
    j    fin
3 error(es); no se escribio ningun archivo.
\end{lstlisting}

Los errores detectados son: registro inexistente o fuera de RV32E, instrucción
desconocida, cantidad incorrecta de operandos, operando vacío, inmediato o salto fuera
de rango, acceso a memoria mal escrito, constante no definida, etiqueta no definida o
repetida, constante repetida, nombre de constante reservado, directiva no soportada y
programa más grande que la RAM de 512 palabras.

Además, el assembler \emph{avisa} (sin fallar) si el programa usa un desplazamiento,
porque el Espino Core lo ejecutaría como una suma y el resultado sería distinto del que
espera quien lee el código.

\subsection{Ejemplo de traducción}

\paragraph{Un branch.} La instrucción \code{bne t0, a3, wrong} de \code{game.s} está en
la dirección \code{0x5C}, y la etiqueta \code{wrong} quedó en \code{0xA8}.

\begin{enumerate}
  \item Desplazamiento: $\texttt{0xA8} - \texttt{0x5C} = \texttt{0x4C} = 76$. En 13
        bits: \code{0\,0000\,0100\,1100}.
  \item Campos: \code{bne} es formato B con opcode \code{1100011} y funct3 \code{001};
        \code{t0} es \code{x5} (\code{00101}) y \code{a3} es \code{x13} (\code{01101}).
  \item Se reparte el desplazamiento: imm[12] = 0, imm[10:5] = \code{000010},
        imm[4:1] = \code{0110}, imm[11] = 0.
\end{enumerate}

\begin{table}[H]
\centering
\small
\begin{tabular}{|c|c|c|c|c|c|c|c|}
\hline
imm[12] & imm[10:5] & rs2 & rs1 & funct3 & imm[4:1] & imm[11] & opcode \\
\hline
0 & 000010 & 01101 & 00101 & 001 & 0110 & 0 & 1100011 \\
\hline
\end{tabular}
\end{table}

Agrupando de a cuatro bits:
\code{0000 0100 1101 0010 1001 0110 0110 0011} = \code{0x04D29663}, que es la línea 24
de \code{game.hex}.

\paragraph{Una constante grande.} \code{li a0, WAIT\_CYCLES}, con
$\text{WAIT\_CYCLES} = 75\,000\,000 = \texttt{0x047868C0}$:

\begin{enumerate}
  \item Los 12 bits bajos son \code{0x8C0}. Su bit 11 es 1, así que con signo valen
        $\texttt{0x8C0} - 4096 = -1856$.
  \item La parte alta se redondea hacia arriba para compensar:
        $(75\,000\,000 + 1856) / 4096 = 18\,311 = \texttt{0x4787}$.
  \item \code{lui a0, 0x4787} carga $\texttt{0x04787000} = 75\,001\,856$, y
        \code{addi a0, a0, -1856} deja $75\,000\,000$.
\end{enumerate}

Codificadas: \code{lui} (formato U, opcode \code{0110111}, rd = \code{01010}) da
\code{0x04787537}, y \code{addi} (formato I, inmediato \code{0x8C0}, rs1 = rd =
\code{01010}, opcode \code{0010011}) da \code{0x8C050513}. Son las líneas 13 y 14 de
\code{game.hex}.
